\documentclass[10pt,a4paper]{article}
\usepackage[utf8]{inputenc}
\usepackage{amsmath}
\usepackage{amsfonts}
\usepackage{comment}
\usepackage{amssymb}
\usepackage{graphicx}
\usepackage{pgfplots}
\pgfplotsset{compat=1.15}
\usepackage{mathrsfs}
\usetikzlibrary{arrows}
\usepackage[ruled,vlined]{algorithm2e}
\usepackage[left=2cm,right=2cm,top=2cm,bottom=2cm]{geometry}
\usepackage{listings}
\usepackage{color}
\definecolor{munsell}{rgb}{0.0, 0.5, 0.69}
\definecolor{dkgreen}{rgb}{0,0.6,0}
\definecolor{gray}{rgb}{0.5,0.5,0.5}
\definecolor{mauve}{rgb}{0.58,0,0.82}
\definecolor{minas}{RGB}{0.244, 0.172, 0.36}
\definecolor{PUJ}{RGB}{44, 86, 151}
\definecolor{PUJ2}{RGB}{43, 93, 156}
\definecolor{PUJ3}{RGB}{20, 52, 107}
\definecolor{cyandk}{rgb}{0.0, 0.72, 0.92}

\author{Iván Andrés Trujillo }

%% code information
\lstset{frame=tb,
  language=Python,
  aboveskip=3mm,
  belowskip=3mm,
  showstringspaces=false,
  columns=flexible,
  basicstyle={\small\ttfamily},
  numbers=none,
  classoffset=1,
  morekeywords={True,False}, keywordstyle=\color{munsell}, 
  classoffset=0, 
  keywordstyle=\color{blue},  
  commentstyle=\color{dkgreen},
  stringstyle=\color{PUJ3},
  numberstyle=\tiny\color{gray},
  breaklines=true,
  breakatwhitespace=true,
  tabsize=4,
}

\begin{document}

\newcommand{\HRule}{\rule{\linewidth}{0.5mm}}

%----------------------------------------------------------------------------------------
%	LOGO SECTION
%----------------------------------------------------------------------------------------

\graphicspath{{/home/ces/Pictures/puj/}}
\includegraphics[width = 4cm]{pujshield.eps}\\[0.5cm] 

\begin{center} 
\textbf{\LARGE Bivariate Analysis using python }\\[0.2cm]
\emph{\LARGE Notes of class}\\[0.3cm] 
\emph{Iván Andrés Trujillo Abella} \\
\textsc{\Large 
}\\[0.2cm] 
\textsc{\large Facultad de ingeniería}\\[0.5cm] 
\HRule \\[0.4cm]
\end{center}
\vspace{1cm}
\section{Analysis}

We need test if two variables have some degree of relation, we need specify the kind of variables involved, categorical or numerical.

\subsection{Exact fisher Test}







\subsection{Contingency table}

\begin{table}[h!]
\begin{tabular}{cccc}
                                      & \textbf{H} & \textbf{M} & \textbf{Total} \\ \cline{2-4} 
\multicolumn{1}{c|}{\textbf{Event A}} & a          & b          & a+b            \\
\multicolumn{1}{c|}{\textbf{Event B}} & c          & d          & c+d            \\
\multicolumn{1}{c|}{\textbf{Total}}   & a+c        & b+d        &               
\end{tabular}
\end{table}

\begin{comment}
yardstick: Criterio
https://latex-tutorial.com/tutorials/tables/ 
    relatedness: parentezco
    digress:divagar
    dull:aburrido
    fashionable: de moda, elegante
    gaze: mirada
    assured:asegurado
    wouldn't (s) no lo haría guedent
    noteworthy:notable nordguordy
    dissapointing: desepcionante
    Nominal variables Crammer's V
    purcell:
    misuse: mal u so
    prevail(s): prevalecer
    inward: interior.
    napkin: servilleta.
    cumbersome: incómodo.
    customary:acostumbrado
    foreshadowing: presagio.
    hatch: incubar
    
    
    
    
 see harvardar probability 19
 
    

\end{comment}


\section{Insights}
There are a relationship between marginal and
joint probabilities 

\subsection{Null hypothesis}

conditional probablities, 

Note if $H_{0}$ it is true then $P(x,y)= P(x)P(y)$ 


the event A underwent $a+b$, note also that 
this ise the sum of marginal frecuencies.
the total of individuals with the feature H is
$a+c$ notice the are also marginal frecuencies.
Note that the probability of a individual underwent the event A is $\frac{a+b}{(a+b)+(c+d)}$ 
now, the expected number that the feature H appear is $ \frac{(a+b)}{(a+b)+(c+d)} (a+c)$.
now we can fill the table with expected values and constrast with the observed values.

If we think in the contigency table as a grid, then we can said that $f_{ij}$ is the frecuency in the row $i$ and $j$ column.

Now we can said $f_{ij}$ to observed frecuencies and $\hat{f}_{ij}$ to the expected frecuencies.

To constract hypothesis 
\begin{equation}
    \chi^{2} = \sum_{i}\sum_{j} \frac{f_{ij}-\hat{f}_{ij}}{\hat{f}_{ij}}
\end{equation}

Degree of fredoms are $(rows-1)(columns-1)$
think in the number of expected observations to eestimate.
\subsubsection{Chi square}


observed - expected values.


Distribution, and degree of freedoms 

$H_{0}$ is independence 
if you get a smaller $p$ value then 
reject this hipothesis



restriction:
expected cell count in each cell is at least five


What is Chocrans Rule?

All expected values are greater than 1
and at least 80\% of expected values
are greater than five.



\subsection{Exact Fisher Test}
One tail p-value.









\subsection{Join, marginal and conditional}

$F(x,y)= P(X<x, Y<y)$ therefore we can 
get $ f(x,y) = \frac{\partial^{2}F}{\partial x \partial y}$


\subsection{Conditional}

 

\subsection{Uniform distribution}
A variable that $X \sim U(a,b)$
means that $ l(b-a) = 1 $ then getting 
$l = \frac{1}{b-a}$
Therefore the pdf of uniform distribution is


\begin{equation}
f(x) =\begin{cases}
\frac{1}{b-a}, & \text{ if $a \leq x   \leq b$} \\
0, & \text{otherwise}
\end{cases}
\end{equation}


think in that $180 ^{\circ} = \frac{\pi}{2} radians $ 
in the unit circle 




\subsection{Unconscious Statistician}
Search the proof of this theorem.




\subsection{Convolution of two formulas}










\subsection{SUM of continuous random variables}

This is important due, Chi squared come from normal distribution:



Here is important to know convolution formula:






\section{Laplace Transform}

\subsection{Potential series}
A potential serie is defined as
$ \sum_{i=0}^{\infty} a_{n}(x-c)^{n}$.

for instance remember that in geometric series $ \sum_{i=0}^{\infty} ax^{i} = \frac{a}{1-x}$ if $-1<x<1$ the last interval as named the \textit{convergence interval.}





\end{document}


